给定二分图 $(L,R,E)$ 及边权 $w_{uv}$，求权值和最大的匹配，允许部分点不匹配。令 $x_{uv}$ 表示是否选择边 $(u,v)$，去掉整数限制得到
\[
\begin{aligned}
\max\quad &\sum_{(u,v)\in E}w_{uv}x_{uv},\\
\text{s.t.}\quad &\sum_{v:(u,v)\in E}x_{uv}\le1\quad(u\in L),\\
&\sum_{u:(u,v)\in E}x_{uv}\le1\quad(v\in R),\\
&x_{uv}\ge0.
\end{aligned}
\]

为左、右两侧每个点的约束分别引入对偶变量 $p_u,q_v$，得到
\[
\begin{aligned}
\min\quad &\sum_{u\in L}p_u+\sum_{v\in R}q_v,\\
\text{s.t.}\quad &p_u+q_v\ge w_{uv}\quad((u,v)\in E),\\
&p_u,q_v\ge0.
\end{aligned}
\]
因此，一组满足约束的点权给出匹配权值和的上界。互补松弛要求：选中的边满足 $p_u+q_v=w_{uv}$，未匹配点的对偶变量为零。若一个匹配和一组对偶可行解满足这些条件，则该匹配最优。边权全为 $1$ 时就是最大匹配。

\textbf{整数性。}若矩阵的任意子式均为 $0$ 或 $\pm1$，称其全幺模。若 $A$ 全幺模且 $b$ 为整数向量，则可行域 $\{x\ge0:Ax\le b\}$ 的所有顶点均为整数点；若有有限最优值，则存在整数最优解。

二分图的点边关联矩阵将一侧点对应的行取负后，成为有向图的关联矩阵，因此全幺模。上面的右端均为 $1$，故存在 $0/1$ 最优解，对应一个匹配。这说明该线性规划松弛不会改变最优值。

\textbf{完美匹配与 KM。}若要求每个点恰好匹配一次，将对应约束改为等式，对偶变量 $p_u,q_v$ 随之变为自由变量，边约束仍为 $p_u+q_v\ge w_{uv}$。将对偶可行解作为顶标，满足等号的边组成相等子图；若其中存在完美匹配，则其权值和等于全部顶标之和，由弱对偶可知最优。

现有二分图最大权匹配模板忽略负权边，补点并将缺失边的权值设为 $0$，求最大权完美匹配后删除其中的零权边，得到允许部分点不匹配的最大权匹配。KM 维护的可行顶标对应这个完美匹配模型的对偶变量；调整顶标并在相等子图中增广，直到得到完美匹配。
